\section*{第 \ref{ch:b3:ode} 章——常微
分方程}

\begin{solution}{exo:b3:ode:1}
(a) 分离变量：$x(t) = \frac1{1 - t}$，定义在
$\intoo{-\infty}1$ 上：在 $T_+ = 1$ 处爆破，且 $x(t) \to
+\infty$，正是 \cref{thm:b3:ode:maximal}(2) 所述。
(b) $x(t) = \tan t$，定义在 $\intoo{-\pi/2}{\pi/2}$ 上，在两端
均爆破。
(c) $x(t) = \frac{1}{1 + \eu^{-t}}$，整体存在：解
始终位于 $\intoo01$ 内，是有界集，因此图像不可能
在有限时间内逃离 $\R\times\R$ 的每个紧集——$T_\pm
= \pm\infty$。注意 (a)、(b) 并不与
\cref{cor:b3:ode:globallinear} 矛盾：$x^2$ 与 $1 + x^2$ 具有
超线性增长。
\end{solution}

\begin{solution}{exo:b3:ode:2}
(a) 这是 \cref{cor:b3:ode:uniqueness} 中当 $t_0
= 0$ 时的估计。该估计是最优的：对于 $x' = Lx$（整体 $L$-利普希茨），两解
之差恰为 $(x_0 - y_0)\eu^{Lt}$。
(b) 该估计说明：时刻 $t$ 的流映射
关于初值是 $\eu^{L\abs t}$-利普希茨的——
即对 $t$ 在紧集上一致连续。对非整体利普希茨的
$f$，在有界集上，
取\omterm{def:b3:groups:action}{轨道}周围的紧管并使用局部常数，重复同样的论证。
\end{solution}

\begin{solution}{exo:b3:ode:3}
(a) $\abs{\sin(tx)} \leq 1$：有界，即线性增长，其中
$\alpha = 0$，$\beta = 1$，由 \cref{cor:b3:ode:globallinear} 在
$U = \R\times\R$ 上得证。
(b) $\abs{tx/(1 + x^2)} \leq \abs t\cdot\frac12$：同样为
次线性增长（事实上在时间紧集上有界）：整体存在。
(c) 记 $X = (x, x')$，$X' = \bigl(\begin{smallmatrix}0 & 1\\ -q(t)
& 0\end{smallmatrix}\bigr)X$：线性，系数连续，
整体存在（\cref{thm:b3:ode:linearstructure} 的
情形）。
(d) 整体存在；用 格朗沃尔 不等式，同 \cref{cor:b3:ode:globallinear}：
$\norm{x(t)} \leq \norm{x(t_0)}\,\eu^{M\abs{t - t_0}}$，其中
$M = \sup\vertiii{A}$。
\end{solution}

\begin{solution}{exo:b3:ode:4}
(a) 对第一个矩阵，$A^2 = -I$：$\eu^{tA} = \cos t\,I + \sin
t\,A = \bigl(\begin{smallmatrix}\cos t & -\sin t\\ \sin t &
\cos t\end{smallmatrix}\bigr)$。若尔当块：$\lambda I$ 与
$N$ 可交换：$\eu^{tA} =
\eu^{\lambda t}\bigl(\begin{smallmatrix}1 & t\\ 0 &
1\end{smallmatrix}\bigr)$。第三个矩阵满足 $A^2 = I$：$\eu^{tA} =
\cosh t\,I + \sinh t\,A$.

(b) 系统为 $X' = \bigl(\begin{smallmatrix}0&1\\-1&0
\end{smallmatrix}\bigr)X + \bigl(\begin{smallmatrix}0\\
\cos\omega t\end{smallmatrix}\bigr)$；借助\omterm{thm:b3:ode:linearstructure}{杜阿梅尔公式}与旋转
的\omterm{thm:b3:ode:linearstructure}{预解式}可得特解：
当 $\omega \neq 1$ 时，$x_p(t) = \frac{\cos(\omega t)}{1 -
\omega^2}$（可直接验证）；对 $\omega = 1$，积分
$\int_0^t\sin(t - s)\cos s\,\dd s = \frac t2\sin t$ 给出
\emph{长期}增长 $x_p = \frac t2\sin t$：共振
——外力在固有频率上注入能量，
振幅线性增长。
\end{solution}

\begin{solution}{exo:b3:ode:5}
(a) $w' = x_1x_2'' - x_1''x_2 = x_1(-px_2' - qx_2) - (-px_1'
- qx_1)x_2 = -p\,w$: $w(t) = w(t_0)\eu^{-\int_{t_0}^tp}$,
永不为零，也不恒等于零。若 $w \neq 0$，则向量
$(x_i, x_i')(t_0)$ 在 $\R^2$ 中线性无关，又由于
解空间维数为 $2$
（由 \cref{thm:b3:ode:linearstructure} 对系统给出），$(x_1,
x_2)$ 为一组基。

(b) 取 $u = \int\frac{\eu^{-\int p}}{x_1^2}$：$x_2 = x_1u$，
$x_2' = x_1'u + \frac{\eu^{-\int p}}{x_1}$，且
\[
x_2'' + px_2' + qx_2
= u\,(x_1'' + px_1' + qx_1) +
\Bigl(-\,p\frac{\eu^{-\int p}}{x_1} +
p\frac{\eu^{-\int p}}{x_1}\Bigr) = 0
\]
（交叉项恰好抵消，展开时须仔细）。对于
$t^2x'' - 2x = 0$，即 $x'' - \frac2{t^2}x = 0$（$p = 0$）
取 $x_1 = t^2$，$u = \int t^{-4} = -\frac1{3t^3}$，故 $x_2
= -\frac1{3t}$，通解为 $at^2 + \frac bt$。
\end{solution}

\begin{solution}{exo:b3:ode:6}
平衡点为 $0, 1$；$f'(x) = 1 - 2x$：$f'(0) = 1 > 0$
（不\omterm{def:b3:ode:stability}{稳定}——邻近解远离，
由显式解可见），$f'(1) = -1 < 0$：渐近\omterm{def:b3:ode:stability}{稳定}
（由 \cref{thm:b3:ode:linearization} 在 $1$ 维情形下得出）。对于
$x(0) \in \intoo01$：此时 $f > 0$，故只要
解始终位于 $\intoo01$ 内，它就递增；它永远不可能到达
$0$ 或 $1$（由唯一性，这两点本身是\omterm{def:b3:groups:action}{轨道}），因此它保持
有界——从而整体存在——并递增趋于极限 $L
\in \intoc{x(0)}1$。若 $f(L) \neq 0$，则 $x' \geq c > 0$
在极限附近，这迫使 $x$ 越过 $L$：故 $f(L) = 0$，$L = 1$；
同理，$x \to 0$（在 $-\infty$ 处）。显式地，$x(t) =
\frac1{1 + C\eu^{-t}}$ 即可验证一切。一般地，
若标量自治方程的解满足 $x'(t_0) = 0$，则 $x(t_0)$
为平衡点，由唯一性知 $x$ 为常数；
否则 $f(x(t))$ 保持定号（它从不为零，
且 $t \mapsto f(x(t))$ 连续）：$x$ 严格
单调——因此非常数的周期解
不可能存在。
\end{solution}

\begin{solution}{exo:b3:ode:7}
(a) $\frac{\dd}{\dd t}H(x,y) = H_xx' + H_yy' = H_xH_y +
H_y(-H_x) = 0$.
(b) $H = \frac{y^2}2 + \frac{x^4}4$ 的水平集为
紧（$H$ 为强制函数），因此解被限制在紧集内：
整体存在且有界（由 \cref{thm:b3:ode:maximal} 得出）。
$(0,0)$ 的\omterm{def:b3:ode:stability}{稳定}性：取 $V = H$，它正定（$H = 0$ 仅在
原点处成立），且 $\dot V = 0$，由 \cref{thm:b3:ode:lyapunov} 即得\omterm{def:b3:ode:stability}{稳定}。
线性化系统 $x' = y$，$y' = 0$ 的矩阵是不可对角化的
幂零矩阵，特征值为 $0$：
\cref{thm:b3:ode:linearization} 不适用（其假设
$\operatorname{Re}\lambda < 0$ 不成立），实际上
线性化系统不\omterm{def:b3:ode:stability}{稳定}（$y_0 \ne 0$ 时发生漂移），而
非线性系统是\omterm{def:b3:ode:stability}{稳定}的：在非双曲平衡点处的线性化
无法说明任何问题。
\end{solution}

\begin{solution}{exo:b3:ode:8}
(a) $\dot V = y\,y' + \sin x\cdot x' = y(-\sin x - cy) +
y\sin x = -cy^2 \leq 0$，且 $V = \frac{y^2}2 + (1 - \cos x)$
它在原点附近的 $\{\abs x < 2\pi\}$ 上正定，
由\cref{thm:b3:ode:lyapunov}可知原点\omterm{def:b3:ode:stability}{稳定}。
(b) 在 $(0,0)$ 处的线性化矩阵为
$\bigl(\begin{smallmatrix}0 & 1\\ -1 &
-c\end{smallmatrix}\bigr)$，其特征多项式为
$\lambda^2 + c\lambda + 1$：根为 $\frac{-c \pm \sqrt{c^2 -
4}}2$——若 $c \geq 2$，两根皆为负实数，
而为复数，其实部 $-\frac c2 < 0$（当 $0 < c < 2$ 时）。
总有 $\operatorname{Re}\lambda < 0$：
由\cref{thm:b3:ode:linearization}可知渐近\omterm{def:b3:ode:stability}{稳定}
（尽管 $\dot V$ 退化）。
(c) 在 $(\pi, 0)$ 处，$\sin(\pi + u) = -\sin u$，线性化为
$\bigl(\begin{smallmatrix}0&1\\1&-c\end{smallmatrix}\bigr)$，
特征多项式为 $\lambda^2 + c\lambda - 1$：两根
异号（$\lambda_+\lambda_- = -1$）。沿不\omterm{def:b3:ode:stability}{稳定}
特征向量，线性系统具有显式的逃逸
解 $\eu^{\lambda_+t}v_+$，其中 $\lambda_+ > 0$，因此
倒立摆对任意阻尼都不\omterm{def:b3:ode:stability}{稳定}。
\end{solution}

\begin{solution}{exo:b3:ode:9}
当 $\alpha \in \intoo01$ 时，除 $x \equiv 0$ 外，每个
\[
x_c(t) = \begin{cases}0 & t \leq c,\\
\bigl((1-\alpha)(t - c)\bigr)^{1/(1-\alpha)} & t \geq c,
\end{cases}
\]
它是 $\mathcal C^1$ 的，且满足方程（指数
$\frac1{1-\alpha} > 1$ 使导数在 $c$ 处为零）：经过
$(0,0)$ 的解构成连续族。当 $\alpha = 1$ 时，$x
\mapsto \abs x$ 是全局 $1$-利普希茨的
（$\abs{\abs a - \abs b} \leq \abs{a - b}$），故
由 \cref{thm:b3:ode:cauchylipschitz} 可知，经过
$0$ 的唯一解为 $x \equiv 0$。分界恰在
$0$ 处的\omterm{def:b3:ode:lipschitz}{局部利普希茨}条件：$\abs x^\alpha$ 在
该点附近差商无界，当 $\alpha < 1$ 时。
\end{solution}

\begin{solution}{exo:b3:ode:10}
(a) 假设 $\norm{x(t_2)} > R$ 对某 $t_2 > 0$ 成立，且
$\norm{x(0)} \leq R$，令 $t_1 = \sup\{t \leq t_2 :
\norm{x(t)} \leq R\}$，则 $\norm{x(t_1)} = R$，且
$\norm{x(t)} > R$ 在 $\intoc{t_1}{t_2}$ 上成立。在该区间上
记 $g(t) = \norm{x(t)}^2$，则 $g'(t) = 2\langle x(t),
F(x(t))\rangle \leq 0$（假设适用：$\norm{x(t)}
\geq R$），于是 $g(t_2) \leq g(t_1) = R^2$，矛盾。
故该球是正向不变的。
(b) 困于紧球内的解不可能有
$T_+ < \infty$（\cref{thm:b3:ode:maximal}(2)）：整体
向前存在。梯度系统：$\frac{\dd}{\dd t}G(x(t)) =
\langle\nabla G, -\nabla G\rangle = -\norm{\nabla G(x(t))}^2
\leq 0$：$G$ 递减，故解保持在 $\{G \leq
G(x_0)\}$ 内，该集合有界（强制性：在大球之外
$G > G(x_0)$），且闭，故紧。逃逸
是不可能的：强制性梯度系统的每一个解都
向前整体存在，永远沿坡下滑。
\end{solution}

\begin{solution}{exo:b3:ode:11}
(a) 比较引理：令 $w = x - y$ 在公共
区间上；$w(0) \geq 0$，且 $w' = x' - y' \geq F(x) - F(y) =
c(t)w$，其中 $c(t) = \frac{F(x) - F(y)}{x - y}$ 在
紧时间区间上有界（$F$ \omterm{def:b3:ode:lipschitz}{局部利普希茨}）；于是
$(w\eu^{-\int c})' \geq 0$，故全程 $w \geq 0$。取
$F(x) = x^2$：$y(t) = \frac1{1 - t}$ 在 $1$ 处爆破，且 $x
\geq y$，只要 $x$ 存在；若 $T_+ > 1$，则 $x$ 将
在 $t = 1$ 处有限，而它却控制着 $y \to \infty$：
荒谬，故 $T_+ \leq 1$。

(b) 反向比较（同一引理，角色互换）：
在 $\intcc01\cap\intco0{T_+}$ 上，由 $t^2 \leq 1$ 得 $x' \leq
x^2 + 1$，而 $z(t) = \tan(t + \frac\pi4)$ 满足 $z' =
z^2 + 1$，$z(0) = 1 = x(0)$，因此在
两者皆有定义期间 $x \leq z$。由于 $z$ 在
$\intco0{\frac\pi4}$ 上有限，$x$ 不可能在
$\frac\pi4$ 之前爆破：$T_+ \geq \frac\pi4$。

(c) 综合得 $\frac\pi4 \leq T_+ \leq 1$（数值上 $T_+
\approx 0.96$）。由此可见：对 $x' = F(t, x)$，若 $F$
关于 $x$ 超线性，解在有限时间内爆破
只要 $\int^\infty\frac{\dd s}{F(s)} < \infty$（
比较解在该有限时间内达到无穷）；
线性增长时积分发散，迫使整体
存在性（见\cref{exo:b3:ode:3}）。这是
与 \cref{pb:b3:complete:1} 中相同的奥斯古德积分，只是此时它决定
的是逃向无穷，而非从零逃离。
\end{solution}

\begin{solution}{exo:b3:ode:12}
(a) $W' = uv'' - u''v = u(-q_2v) - (-q_1u)v = (q_1 -
q_2)\,uv$。

(b) 设 $a < b$ 为 $u$ 的相邻零点，不妨设 $u >
0$ 于 $\intoo ab$ 上（从而 $u'(a) > 0$，$u'(b) < 0$——由唯一性
知其非零，因为 $u(a) = u'(a) = 0$ 将迫使 $u
\equiv 0$）。设 $v$ 在 $\intoo ab$ 内无零点；
不妨设其上 $v > 0$（从而 $v(a), v(b) \geq 0$，由
连续性）。对 (a) 积分：
\[
W(b) - W(a) = \int_a^b(q_1 - q_2)\,uv \;\leq\; 0 .
\]
但 $W(a) = u(a)v'(a) - u'(a)v(a) = -u'(a)v(a) \leq 0$，且
$W(b) = -u'(b)v(b) \geq 0$：故 $W(b) - W(a) \geq 0$。
于是全部取等：$\int(q_1 - q_2)uv = 0$，而 $uv > 0$
在开区间上迫使 $q_1 = q_2$；又 $W(a) =
W(b) = 0$ 迫使 $v(a) = v(b) = 0$；于是 $W \equiv 0$ 在
$\intcc ab$ 上（其导数为零），即\
$(v/u)' = -W/u^2 = 0$，在 $\intoo ab$ 上：$v \propto u$。

（c）取 $q_1 = m$，$u = \sin(\sqrt m(t - t_0))$，其
相邻零点相距 $\pi/\sqrt m$，并取 $q_2 = q \geq
m$：由 (b)，对 $v$ 为 $v'' + qv = 0$ 的任一解，
在每个长为 $\pi/\sqrt m$ 的开区间内都取零（在
退化情形 $q \equiv m$ 中，$v \propto u$ 亦
取零）。若代之以 $q \leq 0$：取 $q_1 = q$，$u =
v$，取 $q_2 = 0$，并取无零点解 $\mathbf 1$，
它是 $v'' = 0$ 的解。若 $v$ 有两个相邻零点，由 (b)
必迫使 $\mathbf 1$ 在二者之间有零点，或退化
情形 $\mathbf 1 \propto v$：二者皆荒谬，$v$ 至多
取零一次。检验：对 $q = 1$，$\sin t$ 每隔 $\pi
= \pi/\sqrt1$，与预言相符；对 $q = -1$，$\sinh t$
恰取零一次，而 $\eu^t$ 从不取零——至多一个
零点，与预言相符。
\end{solution}

\begin{solution}{pb:b3:ode:1}
\textbf{1.} $\dot E = yy' + \sin x\cdot x' = -y\sin x +
y\sin x = 0$：能量为首次积分。在一个\omterm{thm:b3:ode:maximal}{极大
解}上，$y^2 = 2(E_0 + \cos x) \leq 2(E_0 + 1)$，$y$
有界；于是 $\abs{x(t)} \leq \abs{x(0)} +
t\sup\abs y$ 至多线性增长：在任一有限时间
区间内，\omterm{def:b3:groups:action}{轨道}保持在 $\R^2$ 的一个紧集中，故
由\cref{thm:b3:ode:maximal}(2) 可知 $T_\pm = \pm\infty$。
平衡点为 $(k\pi, 0)$；线性化
$\bigl(\begin{smallmatrix}0&1\\
\mp1&0\end{smallmatrix}\bigr)$ 其中 $-\cos(k\pi) = \mp1$：
特征值 $\pm\iu$，对偶数 $k$ 属中心型（
仅凭线性化不能定论），$\pm1$ 对奇数 $k$ 则为鞍点。

\textbf{2.} $E$ 沿解为常数，故每条
\omterm{def:b3:groups:action}{轨道}都限于水平集 $\{y^2 = 2(E_0 + \cos x)\}$。对
$E_0 = -1$：仅有点 $(2k\pi, 0)$。对 $-1 < E_0 <
1$：记 $E_0 = -\cos a$，此集合是互不相交的
闭曲线之并 $y = \pm\sqrt{2(\cos x - \cos a)}$，$x \in
[2k\pi - a, 2k\pi + a]$，每个\omterm{def:b3:ode:stability}{稳定}平衡点周围一条。
对 $E_0 = 1$：曲线 $y = \pm2\cos\frac x2$ 连接
相邻鞍点——即分界线——连同
鞍点本身。对 $E_0 > 1$：两条图像 $y =
\pm\sqrt{2(E_0 + \cos x)}$，对所有 $x$ 有定义，且从不
触及 $y = 0$。

\textbf{3.} $V = E + 1 = \frac{y^2}2 + (1 - \cos x)$
在 $(0,0)$ 处为零，在去心\omterm{def:b3:topology:topology}{邻域}
（$\abs x < 2\pi$）内为正，且 $\dot V = 0 \leq 0$：
由\cref{thm:b3:ode:lyapunov} 得\omterm{def:b3:ode:stability}{稳定}性。不是渐近\omterm{def:b3:ode:stability}{稳定}：
经过 $(a, 0)$（$0 < a$ 充分小）的解
仍留在等值线 $E = -\cos a$ 上，它到原点的距离
为正（该曲线只与 $x$ 轴相交于 $\pm a$）：
$\varphi_t(a, 0) \not\to (0,0)$.

\textbf{4.} 在等值线 $C_a$ 上：无平衡点（$y = 0$
迫使 $x = \pm a$，且 $\sin(\pm a) \neq 0$，对 $0 < a <
\pi$），故速率 $\norm{(y, -\sin x)}$ 具有正的
最小值 $m$（在紧集 $C_a$ 上）。考虑从
$(a, 0)$ 出发的解：在下半平面 $x' = y < 0$，故 $x$
在有限时间内由 $a$ 减至 $-a$（四分之一或半周期
的积分收敛，见第 5 题的分析），到达
$(-a, 0)$；由对称性 $(x, y) \mapsto (x, -y)$，$t
\mapsto -t$ 亦为方程的对称，上半部分沿反向
同样用时 $\frac T2$ 走回：解回到
$(a, 0)$，时刻为 $T$。唯一性
（\cref{cor:b3:ode:uniqueness}）随即推得 $x(t + T) =
x(t)$ 对一切 $t$ 成立，解是周期的。

\textbf{5.} 在 $y > 0$ 的分支上：$y = \sqrt{2(\cos x
- \cos a)}$，又 $\dd t = \frac{\dd x}{y}$；对 $x$
从 $-a$ 积到 $a$ 即得半周期，而对称性 $x
\mapsto -x$ 使积分再减半：
\[
T(a) = 4\int_0^a\frac{\dd x}{\sqrt{2(\cos x - \cos a)}} .
\]
在 $x = a^-$ 处的收敛性：$\cos x - \cos a = \sin(a)(a - x) +
O((a-x)^2)$，且 $\sin a > 0$：被积函数的行为类似于
$\bigl(2\sin a\,(a - x)\bigr)^{-1/2}$，\omterm{def:b3:lebesgue:l1}{可积}。

\textbf{6.} 令 $\sin\frac x2 = k\sin\varphi$，$k =
\sin\frac a2$: $\cos x - \cos a = 2(k^2 - \sin^2\frac x2)
= 2k^2\cos^2\varphi$，且 $\frac12\cos\frac x2\,\dd x =
k\cos\varphi\,\dd\varphi$，于是
\[
T(a) = 4\int_0^{\pi/2}
\frac{\dd\varphi}{\sqrt{1 - k^2\sin^2\varphi}} .
\]
当 $a \to 0^+$ 时 $k \to 0$：对 $k \leq k_0 < 1$，
被积函数受 $(1 - k_0^2\sin^2\varphi)^{-1/2}$ 控制，
在 $\intcc0{\pi/2}$ 上连续：由控制收敛定理得 $T \to
4\cdot\frac\pi2 = 2\pi$。展开 $(1 - u)^{-1/2} = 1 +
\frac u2 + \frac{3u^2}8 + \cdots$ 其中 $u =
k^2\sin^2\varphi$（对 $k < 1$ 正规收敛），
并利用 $\int_0^{\pi/2}\sin^2 = \frac\pi4$：
\[
T(a) = 2\pi\Bigl(1 + \frac{k^2}{4} + O(k^4)\Bigr) :
\]
等时性仅在一阶意义下成立；周期随
振幅增大。

\textbf{7.} 当 $k \uparrow 1$ 时，被积函数递增至 $(1
- \sin^2\varphi)^{-1/2} = \frac1{\cos\varphi}$，其
积分发散：由单调收敛定理，$T(a) \to
+\infty$，当 $a \to \pi^-$ 时。

\textbf{8.} 在分支 $y = 2\cos\frac x2$（$\abs x <
\pi$）：$\frac{\dd x}{2\cos(x/2)} = \dd t$；令 $u = \frac
x2$，$\int\frac{\dd u}{\cos u} =
\ln\tan\bigl(\frac u2 + \frac\pi4\bigr)$，故 $t =
\ln\tan\bigl(\frac x4 + \frac\pi4\bigr)$，即
\[
x(t) = 4\arctan(\eu^t) - \pi,
\qquad y(t) = x'(t) = \frac{4\eu^t}{1 + \eu^{2t}} =
\frac{2}{\cosh t} .
\]
由能量验证：令 $\theta = \arctan\eu^t$，
$\sin2\theta = \frac1{\cosh t}$，故 $\cos x = -\cos4\theta =
-1 + \frac{2}{\cosh^2t}$，且 $E = \frac{y^2}2 - \cos x =
\frac2{\cosh^2t} + 1 - \frac2{\cosh^2t} = 1$：\omterm{def:b3:groups:action}{轨道}
位于分界线上，并对 $y^2 = 2(1 +
\cos x)$ 在 $y > 0$ 时求导即得 $y' = -\sin x$。极限：$x
\to \pm\pi$ 且 $y \to 0$，当 $t \to \pm\infty$ 时。

\textbf{9.} \omterm{def:b3:groups:action}{轨道}向前趋于鞍点 $(\pi, 0)$，
向后趋于 $(-\pi, 0)$，但永不到达：若它到达 $(\pi, 0)$，
设其在有限时刻 $t^*$ 到达，则将有两条不同的\omterm{thm:b3:ode:maximal}{极大解}——
分界线解与位于鞍点的常数解——
都经过同一点
$(t^*, (\pi, 0))$，
与\cref{cor:b3:ode:uniqueness}相矛盾。
鞍点只能渐近逼近。

\textbf{10.} 当 $E_0 > 1$ 时：$y^2 = 2(E_0 + \cos x) \geq
2(E_0 - 1) > 0$：$y$ 保持符号，且 $\abs{x'} = \abs y
\geq \sqrt{2(E_0 - 1)}$：$x$ 严格单调，整体
存在（问题 1），且 $x(t) \to \pm\infty$。由于 $y(t) =
\pm\sqrt{2(E_0 + \cos x(t))}$，而 $\cos$ 以 $2\pi$ 为周期，
$y$ 每当 $x$ 增加 $2\pi$ 时便回到原值；
所需时间为
\[
\tau(E_0) = \int_{x_0}^{x_0 + 2\pi}\frac{\dd
x}{\sqrt{2(E_0 + \cos x)}} =
\int_{-\pi}^{\pi}\frac{\dd x}{\sqrt{2(E_0 + \cos x)}}
\]
（换元；周期性）：$y$ 以 $\tau$ 为周期——
摆以渐近恒定的速率旋转，
该速率为 $2\pi/\tau \approx \sqrt{2E_0}$（对大能量而言）。

\textbf{11.} 方法依次如下：\emph{能量}（$\dot E =
0$）将二维流约化为一维水平曲线；
\emph{有界性}：$y$ 在每条水平曲线上有界，加上不逃逸出紧集，
则给出整体存在性；
\emph{紧性}：闭水平曲线的紧性给出速率界，
进而得到周期性；\emph{唯一性}把首次返回
化为精确的周期性，排除有限时间内到达鞍点，
并区分\omterm{def:b3:groups:action}{轨道}的类型；
\emph{线性化与李雅普诺夫方法}对平衡点分类；
\emph{周期积分}借助
\cref{ch:b3:lebesgue}中的收敛定理作分析。
全过程既不曾拥有——也不需要——一个闭形式的通解：
定性理论提取了运动的每一个特征，
全部来自方程本身。

\textbf{12.} 各部分如下：$W_n = \int_0^{\pi/2}\sin^{2n-1}\varphi
\cdot\sin\varphi\,\dd\varphi = (2n-1)\int_0^{\pi/2}
\sin^{2n-2}\varphi\cos^2\varphi\,\dd\varphi = (2n-1)(W_{n-1}
- W_n)$，因此 $W_n = \frac{2n-1}{2n}W_{n-1}$；又 $W_0 =
\frac\pi2$ 及 $\prod_{j=1}^n\frac{2j-1}{2j} =
\frac{(2n)!}{4^n(n!)^2}$（将 $(2n)! = 2^nn!\prod(2j-1)$ 拆开），
由归纳法即得所示的值。二项级数：$(1 -
u)^{-1/2} = \sum_{n\geq0}c_nu^n$，其中 $c_n =
\frac{(2n)!}{4^n(n!)^2} \in \intoc01$，收敛半径为 $1$。对 $k
\leq k_0 < 1$ 及 $u = k^2\sin^2\varphi$，级数
$\sum c_nk^{2n}\sin^{2n}\varphi$ 关于
$\varphi$ 正规收敛（$c_nk^{2n} \leq k_0^{2n}$），因此
对问题 6 的公式逐项积分是合理的：
\[
T(a) = 4\sum_{n\geq0}c_nk^{2n}W_n
= 4\cdot\frac\pi2\sum_{n\geq0}c_n^2\,k^{2n}
= 2\pi\sum_{n\geq0}
\Bigl(\frac{(2n)!}{4^n(n!)^2}\Bigr)^{\!2}k^{2n} .
\]
取 $c_0 = 1$，$c_1 = \frac12$，$c_2 = \frac38$，则 $T(a) =
2\pi\bigl(1 + \frac{k^2}4 + \frac{9k^4}{64} + O(k^6)\bigr)$，
余项在 $k \leq k_0$ 上一致成立（尾项可由
等比级数控制）。

\textbf{13.} $k = \sin\frac a2 = \frac a2 - \frac{a^3}{48} +
O(a^5)$，于是
\[
k^2 = \frac{a^2}4 - \frac{a^4}{48} + O(a^6),
\qquad
k^4 = \frac{a^4}{16} + O(a^6) ,
\]
且
\[
\frac{T(a)}{2\pi} = 1 + \frac14\Bigl(\frac{a^2}4 -
\frac{a^4}{48}\Bigr) + \frac9{64}\cdot\frac{a^4}{16} +
O(a^6)
= 1 + \frac{a^2}{16} + \frac{11\,a^4}{3072} + O(a^6) ,
\]
因为 $-\frac1{192} + \frac9{1024} = \frac{-16 + 27}{3072}
= \frac{11}{3072}$。

\textbf{14.} 在问题 6 的椭圆形式中，$a \mapsto
k = \sin\frac a2$ 是连续严格递增的
双射，从 $\intoo0\pi$ 到 $\intoo01$，并且对每个
$\varphi \in \intoc0{\pi/2}$，被积函数 $(1 -
k^2\sin^2\varphi)^{-1/2}$ 关于 $k$ 严格递增，于是 $T$
严格递增。连续性：在 $k \leq k_0 < 1$ 上，
被积函数被连续函数 $(1 -
k_0^2\sin^2\varphi)^{-1/2}$ 控制，故
沿 $k' \to k$ 可用控制收敛定理。结合极限 $T \to 2\pi$（当 $a
\to 0^+$ 时（问题 6）及 $T \to +\infty$（当 $a \to \pi^-$
时，问题 7），由严格单调性与介值定理
可知 $T$ 是从 $\intoo0\pi$ 到
$\intoo{2\pi}{+\infty}$ 的双射。

\textbf{15.} 时钟计摆动次数；在振幅趋于零时校准，
则每次摆动记为简谐周期 $2\pi$（
以摆的时间单位计）。在振幅 $a$ 下运行时，
真实周期为 $T(a) = 2\pi\bigl(1 + \frac{a^2}{16} +
O(a^4)\bigr)$：时钟记入 $2\pi$，而真正经过的是 $T(a)$，
故其落后的比例为
\[
\frac{T(a) - 2\pi}{T(a)} = \frac{a^2}{16} + O(a^4) .
\]
当 $a = 0.2$ 弧度（约 $11.5$ 度）时，$\frac{a^2}{16} =
\frac{0.04}{16} = \frac1{400}$，一天有 $86400$ 秒，时钟
每天损失 $86400/400 = 216$ 秒——约三
分半钟——每天如此。因此历史上的两种补救办法为：
强制保持一个极小的恒定振幅（靠擒纵机构实现），或
将约束弯曲，使周期恰好与振幅无关
（惠更斯的摆线颊板，1657 年）。

\textbf{16.} 在 $\tau(E_0) = \int_{-\pi}^{\pi}\frac{\dd
x}{\sqrt{2(E_0 + \cos x)}}$ 对每个固定的
$x$，被积函数关于 $E_0$ 严格递减：$\tau$ 严格
递减。当 $E_0 \downarrow 1$ 时，被积函数递增
逐点趋于 $\bigl(2(1 + \cos x)\bigr)^{-1/2} =
\frac1{2\abs{\cos\frac x2}}$，其在
$\intoo{-\pi}\pi$ 上积分发散（$\cos\frac x2$ 在
$\pm\pi$ 处一阶消失），由单调收敛定理得 $\tau(E_0)
\to +\infty$。当 $E_0 \to +\infty$ 时：
\[
\sqrt{2E_0}\,\tau(E_0) = \int_{-\pi}^{\pi}
\frac{\dd x}{\sqrt{1 + \cos x/E_0}} \longrightarrow 2\pi
\]
由控制收敛定理（对 $E_0 \geq 2$，被积函数
至多为 $\sqrt2$）。故 $\tau \approx 2\pi/\sqrt{2E_0}$：一圈
所需时间等同于以速率 $\sqrt{2E_0}$ 自由旋转的时间，
势能化为一个微小的起伏——与
问题 10 的转速一致。

\textbf{17.} 两坐标轴上有显式解 $t
\mapsto (x_0\eu^{\alpha t}, 0)$ 以及 $t \mapsto (0,
y_0\eu^{-\gamma t})$，以及平衡点
$(0, 0)$：它们都是若干\omterm{def:b3:groups:action}{轨道}的并。向量场为 $\mathcal
C^1$，故\omterm{def:b3:ode:lipschitz}{局部利普希茨}；始于 $Q$ 的一个解
若触及某条坐标轴，则它会经过其中某条\omterm{def:b3:groups:action}{轨道}上的一点，
并由\cref{cor:b3:ode:uniqueness}与该\omterm{def:b3:groups:action}{轨道}重合——
不可能，因为一者位于坐标轴上，另一者则不在。
故 $Q$ 在两个时间方向上均不变。
$Q$ 中的平衡点：由 $x > 0$ 得 $\alpha - \beta y = 0$，
由 $y > 0$ 得 $\delta x - \gamma = 0$：唯一的点为
$(x_*, y_*) = \bigl(\frac\gamma\delta,
\frac\alpha\beta\bigr)$.

\textbf{18.} 沿解，
\[
\dot H = \Bigl(\delta - \frac\gamma x\Bigr)x' +
\Bigl(\beta - \frac\alpha y\Bigr)y'
= (\delta x - \gamma)(\alpha - \beta y) +
(\beta y - \alpha)(\delta x - \gamma) = 0 .
\]
$f(x) = \delta x - \gamma\ln x$ 满足 $f'' = \gamma/x^2 > 0$，
$f'$ 仅在 $x_*$ 处为零，且 $f \to +\infty$ 在
$0^+$ 与 $+\infty$ 两端成立：严格凸且为真函数，极小值为
$f(x_*)$；同理 $g(y) = \beta y - \alpha\ln y$ 的极小值为
$g(y_*)$。故 $H \geq h_*$，等号仅在 $(x_*,
y_*)$ 处成立，且每个下水平集 $\{H \leq h\}\cap Q$ 是紧的：
由 $f(x) \leq h - g(y_*)$，$x$ 被限制在
$\intoo0{+\infty}$ 的一个紧区间内（由真性），$y$ 同理，且该集合
在 $\R^2$ 中是闭的，因 $H \to +\infty$ 于
$Q$ 的边界上。\omterm{thm:b3:ode:maximal}{极大解}停留在其紧的水平集上，
故它不可能在有限时间内离开每一紧集：
由\cref{thm:b3:ode:maximal}可知其为整体的。

\textbf{19.} 取 $h > h_*$，并令 $c = h - g(y_*) >
f(x_*)$。由于 $f$ 自 $+\infty$ 严格递减至
$f(x_*)$（在 $\intoc0{x_*}$ 上），又严格递增
回 $+\infty$（在 $\intco{x_*}{+\infty}$ 上），故方程 $f(x) = c$
恰有两根 $x_- < x_* < x_+$，且 $\{f \leq c\} =
\intcc{x_-}{x_+}$。对 $x \in \intoo{x_-}{x_+}$，$g(y) = h -
f(x) > g(y_*)$ 恰有两根 $y_-(x) < y_* <
y_+(x)$，关于 $x$ 连续（为
$g$ 在
$y_*$ 两侧之严格单调连续限制的逆），且 $y_\pm(x) \to y_*$（当 $x \to x_\pm$ 时）；在 $x =
x_\pm$ 处唯一解为 $y = y_*$。于是 $\{H = h\}\cap
Q$ 为 $y_+$ 与 $y_-$ 的图像在
$\intcc{x_-}{x_+}$ 上的并，在 $(x_\pm, y_*)$ 处粘合为
绕 $(x_*, y_*)$ 的闭曲线——这正是摆的
卵形\omterm{def:b3:groups:action}{轨道}的类比。

\textbf{20.} 令 $C_h = \{H = h\}\cap Q$，$h > h_*$：
$Q$ 的唯一平衡点不在 $C_h$ 上，
故向量场在其上处处不为零。符号为 $x' = \beta x(y_* - y)$，$y' =
\delta y(x - x_*)$：运动在直线 $y
= y_*$ 下方向右，在 $x = x_*$ 右侧向上，上方向左，
左侧向下——为逆时针环绕。沿
下支（开弧）上一点 $(x_0, y_-(x_0))$ 出发的解，
在该支上，此处 $x' > 0$：到达右角点
$B = (x_+, y_*)$ 的时间为
\[
\int_{x_0}^{x_+}\frac{\dd x}{\beta x\,\bigl(y_* -
y_-(x)\bigr)} .
\]
在 $x_+$ 附近取 $x' \in \intoo{x_*}{x_+}$；对 $x \in
\intcc{x'}{x_+}$，$f(x_+) - f(x) \geq f'(x')\,(x_+ - x)$
（$f'$ 在 $x_*$ 之后递增且为正），而
由水平关系及泰勒不等式可得 $g(y_-(x)) -
g(y_*) \leq \frac12\,\bigl(\max g''\bigr)\,(y_* -
y_-(x))^2$，对 $y$ 在 $C_h$ 上取值所成的紧区间成立：故 $y_* -
y_-(x) \geq c\,\sqrt{x_+ - x}$，其中 $c > 0$，而被积函数
为 $O\bigl((x_+ - x)^{-1/2}\bigr)$，故\omterm{def:b3:lebesgue:l1}{可积}——
即问题 5 收敛性的转置。在分支的其余部分，
被积函数连续。故 $B$ 在
有限时间内到达；此处 $y' = \delta y_*(x_+ - x_*) > 0$，\omterm{def:b3:groups:action}{轨道}
进入区域 $x > x_*$，$y > y_*$，攀升至
上角点 $(x_*, y_+)$，由对称估计（
$f$ 与 $g$ 互换角色），并依次绕行四段弧：
经有限时间 $T > 0$ 后，解回到其
出发点。由\cref{cor:b3:ode:uniqueness}，它是
以 $T$ 为周期的——与问题 4 的论证完全相同。

\textbf{21.} 在一条以 $T$ 为周期、位于 $Q$ 中的\omterm{def:b3:groups:action}{轨道}上，$t \mapsto \ln
x(t)$ 是 $\mathcal C^1$ 且以 $T$ 为周期，故
\[
0 = \int_0^T(\ln x)'\,\dd t
= \int_0^T(\alpha - \beta y)\,\dd t
= \alpha T - \beta\int_0^Ty\,\dd t ,
\]
由此得 $\frac1T\int_0^Ty = \frac\alpha\beta$；同理 $0 =
\int_0^T(\ln y)' = \delta\int_0^Tx\,\dd t - \gamma T$ 给出
$\frac1T\int_0^Tx = \frac\gamma\delta$。这些时间平均
恰为平衡点处的值，对每条\omterm{def:b3:groups:action}{轨道}成立，不论振幅如何
——一条无人刻意加入的守恒律。

\textbf{22.} 加入捕捞后，系统仍具有
洛特卡–沃尔泰拉形式，参数为 $\alpha -
\varepsilon$、$\beta$、$\gamma + \varepsilon$、$\delta$（\omterm{def:b3:topology:interior}{内部}平衡点
仍然存在，因为 $\varepsilon <
\alpha$）。对新系统应用第 21 题的结果：
\[
\bar x = \frac{\gamma + \varepsilon}{\delta}
\ \ (\text{猎物平均值上升}),
\qquad
\bar y = \frac{\alpha - \varepsilon}{\beta}
\ \ (\text{捕食者平均值下降}) :
\]
不加选择的捕捞使平衡向
猎物一侧偏移。达·安科纳的数据恰好相反：战争使
捕捞减少，$\varepsilon$ 下降，因此捕食者平均值
$(\alpha - \varepsilon)/\beta$ 上升，猎物平均值
下降——捕获物中鲨鱼所占比例更大，正是
亚得里亚海鱼市场所记录的。这是沃尔泰拉原理，
与农药悖论背后的机制相同：
同时捕杀两个营养级，受益的是被吃掉的那一级。

\textbf{23.} 两次的方法相同：（i）先求第一积分
——摆对应 $E$，此处为 $H$，由分离
$\dd y/\dd x$ 得到——把相平面压缩为曲线族；（ii）
水平集的正常性与紧性，
经由\cref{thm:b3:ode:maximal}给出整体存在性；（iii）
水平集的几何形状——卵形曲线，由 $\cos$ 在该处的形状
及此处 $f$、$g$ 的严格凸性决定——
由积分读出，而非由流读出；（iv）紧卵形线上
不为零的场与\omterm{def:b3:lebesgue:l1}{可积}的角点奇异性
迫使返回时间有限；（v）唯一性
（\cref{cor:b3:ode:uniqueness}）将返回化为
周期性，并排除有限时间内到达平衡点的可能；
（vi）收益——周期展开、平均规律——
来自于对
所得积分应用收敛定理。$x'' = -\sin x$ 与
洛特卡–沃尔泰拉系统都不存在初等闭形式解
（一例为椭圆积分，另一例为超越水平曲线），
而在任何时候都无需它：方程本身
经定性考察，便交出了
整个运动。

\textbf{24.} 在能量水平 $E_0 > 1$ 上，$y^2 = 2(E_0 +
\cos x) > 0$：$y^2$ 的极值在 $\cos x = \pm1$ 处取得，给出
所述的界，在 $x \equiv 0, \pi$ 处取得。在一个周期
$\tau = \tau(E_0)$ 上取时间平均：$x$ 恰增加
$2\pi$，因此
\[
\frac1\tau\int_0^\tau y\,\dd t = \frac{x(\tau) -
x(0)}{\tau} = \frac{2\pi}\tau .
\]
最大与最小速率之比为 $\sqrt{\frac{E_0 + 1}{E_0 -
1}} = 1 + O(E_0^{-1}) \to 1$：高能量下，势能的
$\pm1$ 起伏相对于 $E_0$ 可以忽略，摆
几乎匀速地转动——搓衣板形的势能被抹平。

\textbf{25.} 在 $E_0 \geq 1 + \delta$ 上，
$\tau(E_0) = \int_{-\pi}^{\pi}\frac{\dd x}{\sqrt{2(E_0 +
\cos x)}}$ 的被积函数以 $(2\delta)^{-1/2}$ 为界，且其
关于 $E_0$ 的导数
\[
\partial_{E_0}\frac1{\sqrt{2(E_0 + \cos x)}}
= -\frac{1}{\bigl(2(E_0 + \cos x)\bigr)^{3/2}}
\]
以 $(2\delta)^{-3/2}$ 为界，二者在紧区间
$\intcc{-\pi}\pi$ 上皆\omterm{def:b3:lebesgue:l1}{可积}：积分号下求导
由\cref{thm:b3:lebesgue:paramdiff}可行，于是
得 $\tau'(E_0) < 0$（被积函数严格为负）：
严格递减，且为 $\mathcal C^1$。极限：当 $E_0 \to
\infty$ 时，$\tau \leq \frac{2\pi}{\sqrt{2(E_0 - 1)}} \to 0$；
当 $E_0 \to 1^+$ 时：记 $E_0 + \cos x = (E_0 - 1) +
2\cos^2\frac x2$；被积函数随 $E_0$
减小而增大，故由单调收敛定理
\[
\tau(E_0) \;\nearrow\;
\int_{-\pi}^{\pi}\frac{\dd x}{2\,\abs{\cos\frac x2}}
= +\infty,
\]
极限积分在 $x = \pm\pi$ 处发散（此处
$\abs{\cos\frac x2} \sim \frac{\abs{x \mp \pi}}2$，为
不\omterm{def:b3:lebesgue:l1}{可积}的 $\frac1{\abs\cdot}$）：周期
在趋近分界线时趋于无穷，与第二部分中 $T(a) \to
\infty$ 的结论（从振动一侧）相符。能量
轴的图像为：$E_0 = -1$ 处静止；振动时，$2\pi \nearrow
\infty$，位于 $\intoo{-1}1$ 上；无限慢的分界线在
$E_0 = 1$ 处；转动时，$\infty \searrow 0$，位于其外。
一个积分，摆的整个运动尽在其中。
\end{solution}
